\chapter*{РЕФЕРАТ}
\addcontentsline{toc}{chapter}{РЕФЕРАТ}

РПЗ содержит 65~с., 18~рис., 30~табл., 7~приложений, 23~источника.

Ключевые слова: база данных, спортивная аналитика, NBA, реляционная
модель, \newline PostgreSQL, Redis, кэширование, метрики эффективности, PER, BPM,
REST API, ролевая модель, индексация, производительность.

Объект исследования~--- процессы накопления, обработки и анализа
массивов спортивной статистики NBA.

Предмет исследования~--- структура базы данных
и алгоритмы вычисления комплексных показателей эффективности баскетболистов.

Цель работы~--- разработка базы данных для хранения и анализа
статистики игроков и команд Национальной баскетбольной ассоциации
с реализацией автоматизированных расчётов комплексных показателей эффективности.

В аналитическом разделе проведён сравнительный анализ существующих систем
спортивной статистики, сформулированы функциональные и нефункциональные
требования, формализованы правила расчёта пяти комплексных показателей
(TS\%, eFG\%, USG\%, PER, BPM), выделены восемь сущностей предметной
области, построены ER-диаграмма и диаграмма вариантов использования,
обоснован выбор реляционной модели данных со строковым хранением.

В конструкторском разделе спроектирована реляционная схема из восьми
таблиц, трёх представлений и шести индексов, доказано соответствие 3НФ,
разработаны алгоритмы расчёта метрик и ограничений целостности
(схемы алгоритмов приведены в приложении~А), определена ролевая модель из трёх ролей,
спроектировано REST API из 20~маршрутов и механизм кэширования по паттерну Cache-Aside.

В технологическом разделе обоснован выбор стека (PostgreSQL~15, Redis~7,
FastAPI, React~18), реализована физическая структура базы данных,
разработаны хранимые PL/pgSQL-функции и триггеры (приложение~Г),
приведены примеры SQL-запросов (приложение~В), база данных наполнена
реальной статистикой NBA объёмом более 150\,000 фактов, реализован
пользовательский интерфейс (приложение~Е), проведено регрессионное
тестирование с верификацией по классам эквивалентности.

В исследовательском разделе измерено влияние стратегий индексации,
предвычисленной витрины и кэширования Redis на производительность
аналитических запросов; проведено нагрузочное тестирование
при конкурентном доступе до 100~одновременных подключений;
сформулированы рекомендации по оптимизации.

Разработанный программный комплекс применим в качестве аналитической
платформы для профильных специалистов спортивных организаций.
